%% ============================================================
%%  themes-preview.tex — 主题色预览文档
%% ------------------------------------------------------------
%%  编译后将看到当前主题下所有核心原语的视觉效果。
%%  切换不同主题：修改下方 \settheme{...} 即可。
%% ============================================================

\documentclass[zh]{resume-pro}

%% 可选主题：
%   techblue / businessblack / academicgray / energyorange
%   innovationgreen / professionalpurple / crimson / ocean
\settheme{techblue}

\name{主题预览}
\tagline{resume-pro 视觉规范}
\keywords{展示章节标题 \textbullet{} 卡片 \textbullet{} 标签 \textbullet{} 进度条}
\email{preview@example.com}
\github{hertz}
\location{北京}

\begin{document}

\makeheader

\begin{abstract}
本文档展示 resume-pro 在当前主题下的所有核心视觉组件：标题、列表、技能进度条、技能标签、徽章、项目卡片、开源条目、论文条目、奖项条目等。
\end{abstract}

\sectionTitle{章节标题样式}{\faPaintBrush}

\textbf{加粗文字} / \textit{斜体文字} / \texttt{等宽文字} / \href{https://example.com}{超链接}

\begin{itemize}
  \item 一级列表项使用 \texttt{accentcolor} 渲染的右箭头
  \item 数字与英文混排：99.8\% accuracy on benchmark
  \begin{itemize}
    \item 二级列表项使用小实心圆点
  \end{itemize}
\end{itemize}

\sectionTitle{技能可视化}{\faLaptopCode}

\skillbar{Python \textbullet{} PyTorch}{9}
\skillbar{大模型微调（LoRA / RLHF）}{8}
\skilldots{中文（母语）}{5}
\skilldots{英文（CET-6）}{4}

\vspace{0.4em}
\textbf{技术标签} \quad
\skilltagAccent{LLM} \skilltag{RAG} \skilltag{Agent} \skilltag{LoRA} \skilltag{vLLM}

\vspace{0.3em}
\textbf{徽章} \quad
\badge{Stars}{562} \badge{License}{MIT} \badge{Build}{passing}

\sectionTitle{项目卡片}{\faProjectDiagram}

\projectCard{示例项目 Demo}{2024.01 -- 2024.05}
{一句话说明项目目标与亮点}
{%
\begin{itemize}
  \item 关键技术点 1
  \item 关键技术点 2
\end{itemize}
}
{\faGithub\ github.com/demo/project}

\openSourceItem{example/awesome-llm}{1.2k}{210}{Python}
{开源 LLM 工具库示例项目}
{%
\begin{itemize}
  \item 贡献 PR 若干，修复关键 Bug
\end{itemize}
}

\sectionTitle{论文 / 荣誉}{\faTrophy}

\publicationItem
{An Example Paper on Visual Specification}
{\textbf{Preview Demo}, Co Author}
{NeurIPS 2024 (CCF-A)}
{\faFilePdf\ Paper \quad \faGithub\ Code}

\awardItem{示例荣誉名称}{2024}{说明文字，例如评奖范围、排名、依据等}

\end{document}
